%% ----- develop support ------------------------------------------------------
\usepackage[l2tabu, orthodox]{nag} %% detect and warm about obsolete/outdated/superseded stuff
\usepackage{layout} %% print debug figure of page borders; \layout
\usepackage{lipsum} %% generate dummy text with \lipsum{1-2}
\usepackage{xstring} %% string comparision commands like \IfStrEqCase

%% ----- Umlaute / Symbols ----------------------------------------------------
\usepackage[utf8]{inputenc} %$ erlaubt Sonderzeichen/Umlaute im Quelltext %% ALWAYS ON TOP
\usepackage[T1]{fontenc} %% automatische Trennung von Worten mit Umlauten
%\usepackage{textcomp} %% diverse zeichen wie \texteuro %% TODO REMOVE


%% ----- Language -------------------------------------------------------------
\usepackage[ngerman]{babel} %% set language; new german
%% csquotes loaded at very end of file (preamble): prevent warning


%% ----- Font Family ----------------------------------------------------------
%\usepackage{relsize} %% Schriftgröße relativ festlegen %% TODO REMOVE
%\usepackage[
%    lining %% set numbers on same baseline
%]{nunito} %% uni coporate design
\usepackage{lmodern} %% bessere Font fuer PDFs als latex default
%\usepackage{libertine} \usepackage{libertinust1math} %% eigener fav 1 - text
%\usepackage{mathptmx} %% uni vorgabe - times new romes
%\usepackage[sfdefault]{biolinum}  %% eigener fav 1 - überschriften
%% [sfdefault] only if the base font of the document is to be sans serif

%% ----- fix 'lmodern' not having a bold small capital style -----
\usepackage{fix-cm}
\rmfamily % To load Latin Modern Roman and enable the following NFSS declarations.
% Declare that Latin Modern Roman (lmr) should take
% its bold (b) and bold extended (bx) weight, and small capital (sc) shape,
% from the corresponding Computer Modern Roman (cmr) font, for the T1 font encoding.
\DeclareFontShape{T1}{lmr}{b}{sc}{<->ssub*cmr/bx/sc}{}
\DeclareFontShape{T1}{lmr}{bx}{sc}{<->ssub*cmr/bx/sc}{}


%% ----- Font Style -----------------------------------------------------------
\usepackage{setspace} %% define space betweens lines
\usepackage{geometry} %% define page borders
\usepackage{soul} %% improved text-decoration commands (use \ul isntead of \underline)
%\usepackage{ragged2e} %% improved text alignment; captial version of defaults like \RaggedRight or \begin{FlushLeft} %% TODO incompatible with package 'microtype' option 'spacing' (in style.tex)


%% ----- Table ----------------------------------------------------------------
\PassOptionsToPackage{table}{xcolor} %% include color schema for tables
\usepackage{tabularx} %% automated line-breaks
\usepackage{tabulary} %% control column text direction
\usepackage{booktabs} %% \toprule \midrule[3pt] \cmidrule{1-2} \bottomrule \addlinespace
\usepackage{multicol} %% span column over multiple columns
\usepackage{multirow} %% span row over multiple rows
\usepackage{bigstrut} %% better space alignment in tables; \bigstrut
\usepackage{longtable} %% for long/vertical tables
\usepackage{array} %% extends table column definitions
\usepackage{makecell} %% cellen formatierung mit \makecell
%\usepackage[singlelinecheck=false]{caption} %% do not center single line captions
\usepackage{tablefootnote} %% erlaubt fotenotes in tabellen ohne tricks mit \tablefootnote (NOTE: does hardly work in floating tables)
\newcolumntype{f}[1]{>{\raggedleft\hspace{0pt}}p{#1}} %% Spaltendefinition: rechtsbündig mit definierter Breite
\usepackage{threeparttable} %% handle 'table footnotes' separatly
%\usepackage{footnote} %% provides a 'savenotes' envirnoment, which doesnt rlly work
%\usepackage{footnotehyper} %% same as 'footnote' package
%\makesavenoteenv{tabular}
%\makesavenoteenv{tabulary}

%% ----- Graphics and Media ---------------------------------------------------
\usepackage[dvips,final]{graphicx} %% Einbinden von JPG-Grafiken (in PDF) ermöglichen
\usepackage{graphics} %% \keepaspectratio  %% TODO evtl schon enthalten
\usepackage{wrapfig} %% float/wrap text around figures
%\usepackage{floatflt} %% alternative to wrapfig;
\usepackage{tikz} %% visualisation/drawing
\usetikzlibrary{
    trees, %% directory trees
}
%% ----- Externalize TikZ images -----
%\usetikzlibrary{external}
%\tikzexternalize[prefix=images/,mode=list and make,optimize=false]

%% ----- Math -----------------------------------------------------------------
\usepackage{amssymb} %% fuer /checkmark
\usepackage{amsmath} %% Befehle aus AMSTeX für mathematische Symbole
\usepackage{amsfonts} %% Font fuer amsmath %% TODO evtl auto-loaded by amssymb


%% ----- Listings -------------------------------------------------------------
\usepackage[
    newfloat, %% use 'newfloat' environment
    chapter, %% make listing numbers per chapter
    outputdir={../../out/thesis/} %% move tmp files to output folder
]{minted} %% define with \setminted and print with \inputminted (need latex cmd arg '-shell-escape')

%% or use
%\usepackage{listings} %% define with \lstset and print with \lstinputlisting
%\usepackage{scrhack} %% remove warning from srcreprt due to listings using deprecated '\float@addtolists'

%%% simulatinous usages is not recommended, also may required following fixes:
%% reuse lstlisting counter
%\AtBeginEnvironment{listing}{\setcounter{listing}{\value{lstlisting}}}
%\AtEndEnvironment{listing}{\stepcounter{lstlisting}}


%% ----- Styling --------------------------------------------------------------
\usepackage{xcolor} %% predefined color names
\usepackage{lscape} %% rotates specific page in landscape
\usepackage{caption} %% required for custom listing see '\captionsetup'


%% ----- Sonstiges ------------------------------------------------------------
\usepackage[hyphens]{url} %% URL verlinken, lange URLs umbrechen etc.
\usepackage{chngcntr} %% fortlaufendes Durchnummerieren der Fußnoten
\usepackage{paralist} %% inline aufzählung mit \begin{inparaenum}[a)] (alternative zu enumitem)
\usepackage[inline]{enumitem} %% anpassbare Enumerates/Itemizes
\usepackage{xspace} %% Leerzeichen hinter parameterlosen Commands einfügen
\usepackage{xparse} %% fuer \NewDocumentCommand, ermoeglicht mehrere optionale args %% TODO replace \newcommand
\usepackage{makeidx} %% für Index-Ausgabe (Stichwortverzeichnis) mit \printindex
\usepackage[useregional]{datetime2} %% darstellung eines datum
\usepackage{appendix} %% enable appendix environment


%% ----- Abkürzungsverzeichnis ------------------------------------------------
\ifthenelse{\boolean{isBookPrint}}{
    \PassOptionsToPackage{nohyperlinks}{acronym}
}{}
\usepackage[printonlyused,nolinebreak]{acronym} %% [nolist, printonlyreused, dua, withpage, footnote, smaller]


%% ----- Symbolverzeichnis ----------------------------------------------------
%% TODO does not work
%\usepackage[intoc]{nomencl} %% [intoc] im Inhaltsverzeichnis
%\let\abbrev\nomenclature
%\renewcommand{\nomname}{Symbolverzeichnis}
%\setlength{\nomlabelwidth}{.25\hsize}
%\renewcommand{\nomlabel}[1]{#1 \dotfill}
%\setlength{\nomitemsep}{-\parsep}
%\makenomenclature
%% OR
%\usepackage[intoc, german, noprefix]{nomencl}
%\makenomenclature
%% Mehrere nomencl-Abschnitte:
%% http://www.mrunix.de/forums/showpost.php?p=210422&postcount=29
%\renewcommand{\nomname}{Akronyme}
%\renewcommand{\nompreamble}{\markboth{\nomname}{\nomname}}
%\renewcommand{\nomlabelwidth}{2.5cm}
%\newcommand{\nomaltname}{Symbole}
%\newcommand{\nomaltpreamble}{\markboth{\nomaltname}{\nomaltname}}
%\newcommand{\nomaltpostamble}{}
%\newcommand{\switchnomitem}{S}
%\renewcommand{\nomgroup}[1]{%
%    \ifthenelse{\equal{#1}{\switchnomitem}}{\switchnomalt}{}}
%\newcommand{\switchnomalt}{%
%\end{thenomenclature}
%\renewcommand{\nomname}{\nomaltname}
%\renewcommand{\nompreamble}{\nomaltpreamble}
%\renewcommand{\nompostamble}{\nomaltpostamble}
%\begin{thenomenclature}
%}


%% ----- literaturverzeichnis -------------------------------------------------
\usepackage[
    %backref, %% add '(siehe S. <pages>)' at end of each bib entry, for the usages
    backend=biber, %%is default
    maxnames=20, %% do not shorten author list with 'u.a.'
    style=numeric, %% numeric, alphabetic, authoryear, authortitle, science, ieee
    sorting=none
    %%%% sorting values
    %% nty   := Sort by name, title, year.
    %% nyt   := Sort by name, year, title.
    %% nyvt  := Sort by name, year, volume, title.
    %% anyt  := Sort by alphabetic label, name, year, title.
    %% anyvt := Sort by alphabetic label, name, year, volume, title.
    %% ynt   := Sort by year, name, title.
    %% ydnt  := Sort by year (descending), name, title.
    %% none  := Do not sort at all. All entries are processed in citation order.
]{biblatex} %% http://tug.ctan.org/info/biblatex-cheatsheet/biblatex-cheatsheet.pdf

%% ----- PDF ------------------------------------------------------------------
\usepackage{pdflscape} % pdf-Grafiken quer auf eine Seite
\usepackage{pdfpages} %% import of PDFs into a PDF
\usepackage{hyperref} %% creates links within a document and outside
\usepackage{hypcap} %% ensures hyperref's links to e.x. tables point to that and not the label below


%% ----- END OF FILE (PREAMBLE) -----------------------------------------------
\usepackage[obeyFinal,colorinlistoftodos,textsize=scriptsize]{todonotes} %% in-document todos (must be here ,else graphicx is sad
\usepackage[strict]{csquotes} %% required by babel to ensure qoutes are typeset to language